\section{Análise de Resultados}

\pgfplotstableread[col sep=comma]{dados/regres.csv}\tabelaRegressoes
\pgfplotstablegetelem{0}{m}\of{\tabelaRegressoes} \let\slopeI\pgfplotsretval
\pgfplotstablegetelem{1}{m}\of{\tabelaRegressoes} \let\slopeII\pgfplotsretval
\pgfplotstablegetelem{2}{m}\of{\tabelaRegressoes} \let\slopeIII\pgfplotsretval
\pgfplotstablegetelem{0}{b}\of{\tabelaRegressoes} \let\interceptI\pgfplotsretval
\pgfplotstablegetelem{1}{b}\of{\tabelaRegressoes} \let\interceptII\pgfplotsretval
\pgfplotstablegetelem{2}{b}\of{\tabelaRegressoes} \let\interceptIII\pgfplotsretval


\pgfplotstableread[col sep=comma]{dados/result.csv}\tabelaResult
\pgfplotstablegetelem{0}{Rteo}\of{\tabelaResult} \let\teoI\pgfplotsretval
\pgfplotstablegetelem{1}{Rteo}\of{\tabelaResult} \let\teoII\pgfplotsretval
\pgfplotstablegetelem{2}{Rteo}\of{\tabelaResult} \let\teoIII\pgfplotsretval
\pgfplotstablegetelem{0}{Rexp}\of{\tabelaResult} \let\expI\pgfplotsretval
\pgfplotstablegetelem{1}{Rexp}\of{\tabelaResult} \let\expII\pgfplotsretval
\pgfplotstablegetelem{2}{Rexp}\of{\tabelaResult} \let\expIII\pgfplotsretval
\pgfplotstablegetelem{0}{exat}\of{\tabelaResult} \let\exatI\pgfplotsretval
\pgfplotstablegetelem{1}{exat}\of{\tabelaResult} \let\exatII\pgfplotsretval
\pgfplotstablegetelem{2}{exat}\of{\tabelaResult} \let\exatIII\pgfplotsretval

\pgfplotstablegetelem{2}{Vexp}\of{\tabelaResult} \let\VexpIII\pgfplotsretval
\pgfplotstablegetelem{2}{Vteo}\of{\tabelaResult} \let\VteoIII\pgfplotsretval

\subsection{Montagem 1 - Resistência de Entrada}
\subsubsection{Caso a)}
De acordo com a Lei de Ohm, a tensão e a corrente estão relacionadas de forma linear de forma que, a constante de proporcionalidade é a resistência equivalente do circuito. Assim, representando graficamente a tensão em função da corrente, o declive da reta corresponde ao valor experimental da resistência de entrada do sistema.

Na \cref{gra:mont1-alinA}, encontra-se a regressão linear associada à montagem 1 com C e D ligados, obtida atráves das medições da \cref{tab:mont1-alinA}.

\begin{figure}[!ht]
	\centering
	\begin{tikzpicture}
		\begin{axis}[
				xlabel={Corrente (\unit{\milli\ampere})},
				ylabel={Tensão (\unit{\volt})},
				xmin=0, xmax=42.5, % Adjust these based on your data range
				ymin=0, ymax=17.5,
				grid=both,
				% grid style={dashed, gray!30},
				% legend pos=north west,
				width=0.8\textwidth,
				height=7cm
			]

			\addplot[
				blue,
				mark=*,
				only marks,
			] table [
					col sep=comma,
					x=Current,
					y=Voltage
				] {dados/1a.csv};


			% --- AND PLOT THE REGRESSION LINE ---
			\addplot[red, thick, domain=0:40] {\slopeI * x + \interceptI};

		\end{axis}
	\end{tikzpicture}
	\caption{Gráfico V-I montagem 1 - caso a)}
	\label{gra:mont1-alinA}
\end{figure}

\textbf{Relação entre a tensão e a corrente (C e D ligados):}

\begin{equation}
	V = \pgfmathprintnumber[fixed, precision=4]{\slopeI} \times I
	\label{eq:regLin1}
\end{equation}

Assim, a partir do declive da regressão linear, inferir-se que a resitência de entrada experimental para o circuito com C e D ligados é $R_E=\pgfmathprintnumber[fixed, precision=3]{\expI}\,\unit{\ohm}$.

\paragraph{\textbf{Valor teórico:}}
\paragraph{}

Os valores teóricos foram calculados segundo os valores experimentais da \cref{tab:resist-exper}.

% {\large SAMUEL - Calcular valores do excel e meter aqui}
\begin{equation}
	R_E= R_1+ ({R_2 \parallel R_3})+R_4 = 120,5+\frac{327,0 \cdot 260,0}{327,0 + 260,0}+120,6 = \pgfmathprintnumber[fixed, precision=3]{\teoI}\,\Omega
\end{equation}

De sequência, procedeu-se à determinação da exatidão do valor experimental obtido com a equação \eqref{eq:exatidao}, de forma a avaliar o desvio do resultado experimental em relação ao valor teórico.

% {\large SAMUEL - Calcular valores do excel e meter aqui}
\begin{equation} \label{eq:exatidao}
	\text{Exatidão} = \left(1 - \frac{|R_{\text{exp}} - R_{\text{teo}}|}{R_{\text{teo}}}\right)\times 100 =  \left(1 - \frac{|\pgfmathprintnumber[fixed, precision=3]{\expI} - \pgfmathprintnumber[fixed, precision=3]{\teoI}|}{\pgfmathprintnumber[fixed, precision=3]{\teoI}}\right)\times 100 = \pgfmathprintnumber[fixed, precision=3]{\exatI}\%
\end{equation}


\subsubsection{Caso b)}

Na \cref{gra:mont1-alinB}, encontra-se a regressão linear associada à montagem 2, obtida atráves das medições da \cref{tab:mont1-alinB}.

\begin{figure}[!ht]
	\centering
	\begin{tikzpicture}
		\begin{axis}[
				xlabel={Corrente (\unit{\milli\ampere})},
				ylabel={Tensão (\unit{\volt})},
				xmin=0, xmax=42.5, % Adjust these based on your data range
				ymin=0, ymax=17.5,
				grid=both,
				% grid style={dashed, gray!30},
				% legend pos=north west,
				width=0.8\textwidth,
				height=7cm
			]

			\addplot[
				blue,
				mark=*,
				only marks,
			] table [
					col sep=comma,
					x=Current,
					y=Voltage
				] {dados/1b.csv};


			% --- AND PLOT THE REGRESSION LINE ---
			\addplot[red, thick, domain=0:32.5] {\slopeII * x + \interceptII};

		\end{axis}
	\end{tikzpicture}
	\caption{Gráfico V-I montagem 1 - caso b)}
	\label{gra:mont1-alinB}
\end{figure}

\vspace{10pt}
\textbf{ Relação entre a tensão e a corrente (C
	e D desligados):}
\begin{equation}
	V = \pgfmathprintnumber[fixed, precision=3]{\slopeII} \times I
	\label{eq:regLin2}
\end{equation}


Da mesma forma, a partir do declive da regressão linear, inferir-se que a resitência de entrada experimental para o circuito com C e D desligados é $R_E= \pgfmathprintnumber[fixed, precision=3]{\expII} \,\unit{\ohm}$.

\paragraph{\textbf{Valor teórico:}}

\begin{equation*}
	R_E= R_1+ R_3 +R_4 = 120,5+ 260 +120,6 = \pgfmathprintnumber[fixed, precision=3]{\teoII}\,\Omega
\end{equation*}

Posteriormente, determinou-se a exatidão do valor experimental obtido através da equação \eqref{eq:exatidao}.

\begin{equation*}
	\text{Exatidão} = \left(1 - \frac{|\pgfmathprintnumber[fixed, precision=3]{\expII}-\pgfmathprintnumber[fixed, precision=3]{\teoII}|}{\pgfmathprintnumber[fixed, precision=3]{\teoII}}\right)\times 100 = \pgfmathprintnumber[fixed, precision=3]{\exatII}\%
\end{equation*}


\subsection{Montagem 2 - Fonte real}
Na \cref{gra:mont2}, encontra-se a regressão linear associada à montagem 2, obtida atráves das medições da \cref{tab:mont2}.

\enlargethispage{2\baselineskip}

\begin{figure}[!ht]
	\centering
	\begin{tikzpicture}
		\begin{axis}[
				xlabel={Corrente (\unit{\milli\ampere})},
				ylabel={Tensão (\unit{\volt})},
				xmin=0, xmax=18, % Adjust these based on your data range
				ymin=0, ymax=8,
				grid=both,
				% grid style={dashed, gray!30},
				% legend pos=north west,
				width=0.8\textwidth,
				height=7cm
			]

			\addplot[
				blue,
				mark=*,
				only marks,
			] table [
					col sep=comma,
					x=Current,
					y=Voltage
				] {dados/2.csv};


			% --- AND PLOT THE REGRESSION LINE ---
			\addplot[red, thick, domain=0:18] {\slopeIII * x + \interceptIII};


		\end{axis}
	\end{tikzpicture}
	\caption{Gráfico V-I montagem 2}
	\label{gra:mont2}
\end{figure}

Obteve-se assim a seguinte equação característica da fonte de tensão:
\begin{equation*}
	V = \pgfmathprintnumber[fixed, precision=3]{\slopeIII} \times I + \pgfmathprintnumber[fixed, precision=3]{\interceptIII}
	\label{eq:regLin3}
\end{equation*}

A equação da fonte real é dada por: V = $-R_i$I $+ V_0$. Dessa forma, obtém-se $V_0 = \pgfmathprintnumber[fixed, precision=3]{\VexpIII}\,\unit{V}$ e $R_i = \pgfmathprintnumber[fixed, precision=3]{\expIII}\,\unit{\ohm}$.


\paragraph{\textbf{Valor teórico:}}

\vspace{10pt}
Anulando a fonte (curto-circuito):

\begin{align*}
	               & R_{eq} = R_2 + (R_3 \parallel R_1 + R_4)                                                                              \\
	               & R_1 + R_4 = 241,1 \, \Omega                                                                                           \\
	\Rightarrow \, & R_{eq} = 327 + \frac{260 \cdot 241,1}{260 + 241,1} \approx \pgfmathprintnumber[fixed, precision=3]{\teoIII} \, \Omega
\end{align*}

Em circuito aberto entre 1 e 2, não circula corrente em $R_2$, logo não existe queda de tensão nessa resistência. A queda de tensão a calcular é a queda de tensão em $R_3$.

\begin{align*}
	V_{Th} = V_3 = V_s \, \frac{R_3}{R_1+R_3+R_4} = 15 \cdot \frac{260}{501} = \pgfmathprintnumber[fixed, precision=3]{\VteoIII}\,\unit{V}
\end{align*}

Por fim, calculou-se a exatidão do valor experimental obtido com a equação \eqref{eq:exatidao}.
\begin{equation*}
	\text{Exatidão} = \left(1 - \frac{|448 - 452|}{452}\right)\times 100 = \pgfmathprintnumber[fixed, precision=3]{\exatIII}\%
\end{equation*}
